\section{Three traps that would each have produced a positive result}
\label{sec:traps}

\placefig{fig-traps}

\paragraph{Trap 1: the direction measures position.} False completions come later in an episode than true ones, so anything that grows with the episode separates them. A first difference-of-means direction between \fsp{} and \fsm{} steps (\prvPos{} and \prvNeg{} steps), cross-validated by task, peaked at AUROC \prvAucPost{} against \prvBaseStep{} for the step index and \prvBaseCtx{} for context length. Reported without baselines, \prvAucPost{} reads as a solid detector; in fact it was worse than counting steps. The official analysis, run once \fsp{} steps passed the pre-registered 40 (\ofPos{} \fsp, \ofNeg{} \fsm, \ofTasks{} tasks, both passes; exploratory because the second pass stopped early), confirms it (Table~\ref{tab:fs-auc}, Fig.~\ref{fig:traps}b). The best direction, at layer \ofResidLayer{} on the last prefill token, reaches \ofResidAuc{} (logistic probe \ofResidLr) against \ofBaseStep{} and \ofBaseCtx{} for the two position baselines, and the token before the action does worse (\ofPostAuc). Controlling for position directly, a logistic model on step and log context length scores \ofPosOnly; adding the best direction raises it by at most \ofPosGain, at a layer chosen after the fact, within noise. We find no false-success signal in the residual stream beyond the length of the episode, which is consistent with the failure of every direction-based intervention below. A value-axis direction in the style of \citet{jiang2026value} shows the same confound in a separate construction (appendix~\ref{app:value}). Our protocol therefore matches positives and negatives on step and context length before fitting, and reports both baselines next to every AUROC.

\placefig{tab-fs-auc}

\paragraph{Trap 2: regression to the mean.} Agent outcomes are stochastic at $T=0.8$. Selecting tasks because they failed in one run selects partly for bad luck. Re-running the unmodified agent on the same tasks, only \rtmFsAgain/\rtmFsN{} (\rtmFsAgainPct\%, 95\% CI \rtmFsAgainLo--\rtmFsAgainHi) pass-1 false successes are false successes again (an earlier snapshot of the second pass, on fewer tasks, gave 42\%), and \rtmFsSucc{} of them now succeed (Fig.~\ref{fig:traps}c). % source of the 42%: data/local/reports/CONTEXT_2026-11-07.md "Régression vers la moyenne"
A patch evaluated as ``false-success tasks before versus after'' would therefore be credited with about \rtmFakeEffect{} points of reduction while doing nothing. Our validation compares two fresh runs, original and patched, on a frozen reserve of \texttt{confirm\_tasks.json} tasks that no intervention was built on (\S\ref{sec:protocol}).

\paragraph{Trap 3: detecting is not causing.} The 0.8B loop direction is a textbook detection result (Fig.~\ref{fig:traps}a): task-grouped AUROC \mwAucBestExact{} at layer \mwAucBestLayer, against \mwAucStep{} for the step index and \mwAucCtx{} for context length, replicated at \mwAstraAuc{} on an independent refit. It fails every causal test we ran (\S\ref{sec:mw}). Projecting it out of every writer erases the marker from the residual stream (the loop/non-loop projection gap at layer 10 falls from 0.1653 to $-0.00008$) and leaves loop actions as likely as before relative to healthy ones. % source: data/pod_astra/REPORT.md "Contrôle mécanistique : effacement global du même u10"
The marker can be deleted without touching the behaviour it marks.

\subsection{The protocol these traps impose}
\label{sec:protocol}
\begin{enumerate}
\item \textbf{Free generation is the only efficacy measure.} Teacher-forced log-probabilities of a recorded action mix format, arguments and end-of-sequence, and on the 4B the decision is taken during reasoning, after which $P(\texttt{answer})$ is close to 1 for true and false completions alike. % source: ~/.claude/skills/weight-surgery/SKILL.md "Defect-specific notes" [evidence: Holo-3.1-4B]
Each candidate regenerates the full reasoning ($k=4$ per context, fixed seeds shared across arms, production sampling with \texttt{top\_k=20}), and invalid or truncated outputs stay in every denominator.
\item \textbf{A gate fixed before any data.} On \fsp{} contexts of tasks not used to build the intervention, $P(\texttt{answer})$ must fall by at least 25 points and 50\% relative. The \fsm{} answer rate and the rate of valid non-answer actions on mid-episode contexts may each fall by at most 5 points, and invalid outputs may rise by at most 5 points. % source: data/pod_astra/4b/PROTOCOL.md section C
The \fsm{} criterion is what separates ``less false success'' from ``less success''.
\item \textbf{Controls.} A random direction with the same norm and scope, an identity hook that measures the noise floor, and a fixed budget of generation batches per hypothesis, counted before execution, failures included. Two consecutive settings without progress stop a hypothesis.
\item \textbf{Disjoint confirmation.} The official direction analysis is gated at 40 admissible \fsp{} steps. The \texttt{confirm\_tasks.json} reserve is excluded from every construction and screening, and was to be read only to confirm a committed candidate, with both arms re-run fresh. In the event no candidate was admitted, and five non-admitted arms were run on it (\S\ref{sec:online}); we list this as a limitation (\S\ref{sec:limits}).
\end{enumerate}
